Introduction
Phonology
Morphology
Derivational Morphology
Syntax
Semantic Fields and Pragmatics
Writing System
Examples
Lexicon



Intro
Phonology
Vowels
Consonants
Stress
Writing System
Grammar
Nouns (and derivational morphology)
** Cases, pre-/post-positions and such. Everything related to nouns.
Verbs (and derivational morphology)
** Moods, tenses, aspects, voices, sequence of tenses (tense harmony) and conditional sentences...
Adjectives (or stative, attributive, adjectival verbs)
Adverbs (or words modifying verbs and adjectives)
Anything else depending on the conlang
Number system and counter words, etc...
Pragmatics, register levels and such...



1.	Introduction
	1.1.	Grammar overview
2.	Orthography
	2.1.	Alphabet
	2.2.	Braille
3.	Phonology
	3.1.	Consonants
	3.2.	Vowels
	3.3.	Diphthongs
4.	Phonotactics
	4.1.	Syllable structures
	4.2.	Length
	4.3.	Stress
5.	Grammar
	5.1.	Sentence structure
	5.2.	Word classification
	5.3.	Numbers
		5.3.1.	Cardinal numbers
		5.3.2.	Ordinal numbers
		5.3.3.	Nominal numbers
	5.4.	Nouns and postpositions
		5.4.1.	Number
		5.4.2.	Postpositions and case markers
			5.4.2.1.	Core cases
			5.4.2.2.	Postpositions
		5.4.3.	Nominal affixes
	5.5.	Pronouns
		5.5.1.	Personal pronouns
		5.5.2.	Demonstrative and quantifying determiners
		5.5.3.	Extended quantifier pronouns and adjectives
		5.5.4.	Reciprocal pronouns
		5.5.5.	Reflexive pronouns
	5.6.	Verbs and modal particles
		5.6.1.	Declensions
		5.6.2.	Verbal lexical modification
			5.6.2.1.	Prefixes: Separable
			5.6.2.2.	Prefixes: Inseparable
			5.6.2.3.	Suffixes
		5.6.3.	Moods
		5.6.4.	Voices
		5.6.5.	Verbal anaphor
	5.7.	Adjectives
		5.7.1.	Comparative and superlative formation
		5.7.2.	Use as verbs
	5.8.	Adverbs
		5.8.1.	Comparative and superlative formation
		5.8.2.	Temporality
	5.9.	Conjunctions
		5.9.1.	Coordinating conjunctions
		5.9.2.	Subordinating conjunctions
		5.9.3.	Correlative conjunctions
	5.10.	Particles
6.	Basic phrases
7.	Appendices
	7.1.	Abbreviations
	7.2.	Terminology
	7.3.	Selected translations
 


1. Introduction
2. Phonology
  2.1. Consonants
    2.1.1. Noteworthy Consonantal Allophones
  2.2. Vowels
    2.2.1. Basic Vowel Qualities
    2.2.2. Diphthongs and Triphthongs
    2.2.3. Phonation
    2.2.4. Full Vowel Listing
      2.2.4.1. Vowels Listed by Phonation Type
      2.2.4.2. Traditional Listing of Vowels
  2.3 Note on Romanization
  2.4. Phonotactics
      2.4.1. Onsets
      2.4.2. Nuclei
      2.4.3. Codas
  2.5. Stress
      2.5.1. Lexical Stress
      2.5.2. Prosodic Stress
      2.5.2. Stressed Vowel Pitch
3. Morphology
  3.1. Nouns
    3.1.1 Plurals
      3.1.1.1. Associative Plurals
    3.1.2. Pronouns
      3.1.2.1. Personal Pronouns
      3.1.2.2. Interrogative Pronouns
      3.1.2.3. Third Person Pronouns as Topic Markers
      3.1.2.4. Relative Pronouns
      3.1.2.5. Archaic or Dialectal Pronouns
    3.1.3. Demonstrative Constructions
      3.1.3.1. Locative Pronouns
      3.1.3.2. Demonstrative Adjectives
      3.1.3.2. Demonstrative Pronouns
  3.2. Verbs
    3.2.1. Dynamic Verbs
    3.2.2. Perfective Verbs
    3.2.3. Stative Verbs
    3.2.4. Semantic Similarities Between Verbs of Different Aspectual Classes
    3.2.5. The Copula
    3.2.6. Auxiliary Verbs
      3.2.6.1. Alignment Auxiliaries
      3.2.6.2. Perfect Auxiliaries
      3.2.6.3. Imperfect Auxiliaries
      3.2.6.4. Modal Auxiliaries
      3.2.6.5. Order of Auxiliaries	
    3.2.8. Transitivity
      3.2.8.1. Explicit Marking of Intransitive Verbs
      3.2.8.2. Explicit Marking of Transitive Verbs
    3.2.9. Verb Nominalization
    3.2.7. Pseudo-Prepositional Verbs
      3.2.7.??. (not yet organized)
  3.3. Conjunctions
    3.3.??. (not yet organized)
  3.4. Interjections and Sentence Final Particles
    3.4.1. Particles of Request
    3.4.2. Particles of Suggestion
    3.4.3. Particles of Affirmation
    3.4.4. Verb-Particle Relationships
4. Derivational Morphology
  1. Noun-Noun Compounds
  2. Noun-Verb Compounds
  3. Verb-Noun Compounds
  4. Serial Verb Compounds
  5. Reduplication
5. Syntax
1. Standard Word Order
  1.1. Topic vs. Subject
    1.1.1. Topic Fronting
2. Question forms
  2.1 Content Questions
  2.2 Polar Questions
3. Comparatives & Superlatives
4. Counting
  4.??. (not organized)
7.  Writing System
1. Fkeuswa
  1.1. Semantic Primitives
  1.2. Semantic Radicals
  1.3. Phonetic Complements
  1.3.1. Onset Complements
  1.3.2. Onset-Nucleus Complements
  1.3.3. Morpheme Complements
  1.4 Simplification
  1.5. Typography
    1.5.1. Punctuation Marks
    1.5.2. Standard Paragraphing
3. Fkeumgerswa
  3.1. Consonant Signs
  3.2. Vowel Signs
    3.2.1. Determination of Vowel Sound
  3.3. Consonant Subscripts
  3.4. Punctuation





2 Phonology
  2.1 Phoneme Chart
  2.2 Consonants
  2.3 Vowels
  2.4 Minimal pairs and examples
  2.5 Phonological Processes
    2.5.1 Nasal place assimilation
    2.5.2 Gemination
    2.5.3 Copy vowels
    2.5.4 Word-initial fortition
  2.6 Phonotactic Restrictions
  2.7 Stress

3 Nouns
  3.1 Noun Classes
    3.1.1 Terrestrial
    3.1.2 Aquatic
    3.1.3 Synthetic
    3.1.4 Phenomenal
  3.2 Noun Cases
    3.2.1 Nominative
    3.2.2 Accusative
    3.2.3 Genitive
    3.2.4 Dative
    3.2.5 Locative
  3.3 Noun Number

4 Verbs
  4.1 Slot 1
  4.2 Slot 2
  4.3 Slot 3
  4.4 Slot
  4.5 Slot
  4.6 Slot
  4.7 Voice

5 Adjectives and Adverbs

6 Postpositions

7 Pronouns and Determiners

8 Comparatives

9 Numerals

10 Derivations

11 Syntax
  11.1 Noun phrases
  11.2 Serial verbs
  11.3 Question formation

12 Writing System

13 Sample Translations